\documentclass[11pt,a4paper]{article}

% Pakete
\usepackage[utf8]{inputenc}
\usepackage[T1]{fontenc}
\usepackage[ngerman]{babel}
\usepackage[left=18mm,right=18mm,top=18mm,bottom=18mm]{geometry}
\usepackage{helvet}
\renewcommand{\familydefault}{\sfdefault}
\usepackage{amsmath,amssymb}
\usepackage{graphicx}
\usepackage{tikz}
\usetikzlibrary{arrows.meta,positioning,patterns}
\usepackage{tcolorbox}
\tcbuselibrary{skins,breakable}
\usepackage{enumitem}
\usepackage{tabularx}
\usepackage{colortbl}
\usepackage{qrcode}
\usepackage{xcolor}
\usepackage[hidelinks]{hyperref}

% Fachfarbe Physik
\definecolor{physik}{HTML}{2c5aa0}
\definecolor{physiklight}{HTML}{e8f0fa}

% Box-Stile
\tcbset{
    merkbox/.style={
        colback=physiklight,
        colframe=physik,
        fonttitle=\bfseries,
        title=#1,
        boxrule=1pt,
        arc=2pt,
        left=5pt,
        right=5pt,
        top=5pt,
        bottom=5pt
    },
    formelbox/.style={
        colback=white,
        colframe=physik,
        boxrule=1pt,
        arc=2pt,
        left=6pt,
        right=6pt,
        top=6pt,
        bottom=6pt
    }
}

% Kopfzeile
\usepackage{fancyhdr}
\pagestyle{fancy}
\fancyhf{}
\lhead{\small Physik 12 GK}
\chead{\small Übung: Testvorbereitung}
\rhead{\small Name: \underline{\hspace{4cm}}}
\renewcommand{\headrulewidth}{0.4pt}

% Keine Einrückung
\setlength{\parindent}{0pt}
\setlength{\parskip}{6pt}

\begin{document}

% Titel
\begin{minipage}[t]{0.65\textwidth}
\vspace{0pt}
{\Large\textbf{Übung: Photoeffekt, De-Broglie, Doppelspalt}}\\[0.3em]
{\normalsize Fr 23.01.2026 — Testvorbereitung (Test: Mi 28.01.)}
\end{minipage}
\hfill
\begin{minipage}[t]{0.32\textwidth}
\vspace{0pt}
\raggedleft
\begin{tabular}{@{}c@{\hspace{8pt}}c@{}}
\href{https://devbox42.github.io/sciencesim/physik/kl12-GK/01-photoeffekt/05-DS-uebung/SIM-05-elektronenbeugung.html}{\qrcode[height=1.1cm]{https://devbox42.github.io/sciencesim/physik/kl12-GK/01-photoeffekt/05-DS-uebung/SIM-05-elektronenbeugung.html}} &
\href{https://devbox42.github.io/sciencesim/physik/kl12-GK/01-photoeffekt/05-DS-uebung/LP-05-uebung.html}{\qrcode[height=1.1cm]{https://devbox42.github.io/sciencesim/physik/kl12-GK/01-photoeffekt/05-DS-uebung/LP-05-uebung.html}} \\
{\scriptsize e-Beugung} & {\scriptsize Lernpfad}
\end{tabular}
\end{minipage}

\vspace{0.3em}
\hrule height 1pt
\vspace{0.5em}

% ============================================================
% FORMELSAMMLUNG
% ============================================================

\begin{tcolorbox}[formelbox]
\textbf{Formelsammlung (für den Test erlaubt)}

\vspace{0.3em}
\begin{tabular}{@{}p{0.48\textwidth}p{0.48\textwidth}@{}}
\textbf{Photoeffekt:} & \textbf{De-Broglie:} \\[0.2em]
$E_{\text{Photon}} = h \cdot f$ & $\lambda = \dfrac{h}{p} = \dfrac{h}{m \cdot v}$ \\[0.8em]
$E_{\text{kin,max}} = h \cdot f - W_A$ & $p = \dfrac{h}{\lambda}$ \\[0.8em]
$E_{\text{kin}} = e \cdot U_G$ & $E_{\text{kin}} = \dfrac{p^2}{2m}$ \\
\end{tabular}

\vspace{0.3em}
\textbf{Konstanten:} $h = 6{,}63 \cdot 10^{-34}$ J$\cdot$s $= 4{,}14 \cdot 10^{-15}$ eV$\cdot$s \quad | \quad $e = 1{,}60 \cdot 10^{-19}$ C \quad | \quad $m_e = 9{,}11 \cdot 10^{-31}$ kg
\end{tcolorbox}

\vspace{0.5em}

% ============================================================
% TEIL A: PHOTOEFFEKT
% ============================================================

{\large\textbf{Teil A: Photoeffekt}} \hfill \textbf{(12 BE)}

\vspace{0.3em}

\textbf{Aufgabe 1:} Eine Kaliumkathode ($W_A = 2{,}3$ eV) wird mit Licht der Wellenlänge $\lambda = 400$ nm bestrahlt.

\begin{enumerate}[label=\alph*)]
\item Berechne die Energie eines Photons in eV. \hfill (3 BE)

\vspace{2cm}

\item Berechne die maximale kinetische Energie der ausgelösten Elektronen. \hfill (2 BE)

\vspace{1.5cm}

\item Bestimme die Grenzwellenlänge für Kalium. \hfill (3 BE)

\vspace{2cm}

\end{enumerate}

\textbf{Aufgabe 2:} Das Diagramm zeigt die Einsteinsche Gerade für zwei Metalle A und B.

\begin{minipage}[t]{0.5\textwidth}
\vspace{0pt}
\begin{tikzpicture}[scale=0.8]
    \draw[gray!30, thin, xstep=1, ystep=0.5] (0,-2) grid (7,2.5);
    \draw[->, thick] (0,-2) -- (0,2.8) node[above] {$E_{\text{kin}}$ (eV)};
    \draw[->, thick] (0,0) -- (7.3,0) node[right] {$f$ ($10^{14}$ Hz)};

    % Achsenbeschriftung
    \foreach \x in {2,4,6} \node[below] at (\x,0) {\small \x};
    \foreach \y in {-1,1,2} \node[left] at (0,\y) {\small \y};

    % Gerade A (W_A = 1.5 eV, f_G = 3.6)
    \draw[physik, thick] (3.6,0) -- (7,1.4);
    \node[physik, right] at (6.5,1.2) {A};

    % Gerade B (W_A = 2.5 eV, f_G = 6)
    \draw[red!70!black, thick] (6,0) -- (7,0.4);
    \draw[red!70!black, thick, dashed] (0,-2.5) -- (6,0);
    \node[red!70!black, right] at (6.5,0.2) {B};
\end{tikzpicture}
\end{minipage}
\hfill
\begin{minipage}[t]{0.45\textwidth}
\vspace{0pt}
\begin{enumerate}[label=\alph*), start=4]
\item Welches Metall hat die größere Austrittsarbeit? Begründe. \hfill (2 BE)

\vspace{1.8cm}

\item Lies die Grenzfrequenz von Metall A ab und berechne $W_A$. \hfill (2 BE)

\vspace{1.5cm}

\end{enumerate}
\end{minipage}

\newpage

% ============================================================
% TEIL B: DE-BROGLIE-WELLENLÄNGE
% ============================================================

{\large\textbf{Teil B: De-Broglie-Wellenlänge}} \hfill \textbf{(10 BE)}

\vspace{0.3em}

\textbf{Aufgabe 3:} Elektronen werden mit einer Spannung $U_B = 5{,}0$ kV beschleunigt.

\begin{enumerate}[label=\alph*)]
\item Berechne die kinetische Energie der Elektronen in Joule. \hfill (2 BE)

\vspace{1.5cm}

\item Berechne den Impuls der Elektronen. \hfill (3 BE)

\textit{Hinweis:} $E_{\text{kin}} = \dfrac{p^2}{2m}$ $\Rightarrow$ $p = \sqrt{2 \cdot m_e \cdot E_{\text{kin}}}$

\vspace{2cm}

\item Berechne die De-Broglie-Wellenlänge der Elektronen in pm. \hfill (2 BE)

\vspace{1.5cm}

\end{enumerate}

\textbf{Aufgabe 4:} Ein Elektron hat die De-Broglie-Wellenlänge $\lambda = 100$ pm.

\begin{enumerate}[label=\alph*), start=4]
\item Berechne den Impuls des Elektrons. \hfill (2 BE)

\vspace{1.5cm}

\item Vergleiche: Welche Wellenlänge hat ein Photon mit gleichem Impuls? \hfill (1 BE)

\vspace{1cm}

\end{enumerate}

\vspace{0.5em}

% ============================================================
% TEIL C: DOPPELSPALT UND KOMPLEMENTARITÄT
% ============================================================

{\large\textbf{Teil C: Doppelspalt und Komplementarität}} \hfill \textbf{(8 BE)}

\vspace{0.3em}

\textbf{Aufgabe 5:} Elektronen werden einzeln durch einen Doppelspalt geschickt.

\begin{enumerate}[label=\alph*)]
\item Beschreibe, was man auf dem Schirm beobachtet, wenn nur wenige Elektronen auftreffen. \hfill (2 BE)

\vspace{1.5cm}

\item Beschreibe, was man beobachtet, wenn sehr viele Elektronen aufgetroffen sind. \hfill (2 BE)

\vspace{1.5cm}

\item Erkläre, warum dieses Verhalten mit dem klassischen Teilchenbild nicht vereinbar ist. \hfill (2 BE)

\vspace{1.5cm}

\item Was passiert mit dem Muster, wenn man einen Detektor aufstellt, der misst, durch welchen Spalt jedes Elektron geht? Begründe mit dem Komplementaritätsprinzip. \hfill (2 BE)

\vspace{1.5cm}

\end{enumerate}

\begin{center}
\rule{10cm}{0.5pt}

\textbf{Gesamt:} \underline{\hspace{1cm}} / 30 BE
\end{center}

\end{document}
